\documentclass[12pt, a4paper]{article}

% Required Packages
\usepackage{amsmath, amssymb, graphicx}
\usepackage{kotex} % For Korean text
\usepackage{geometry}
\usepackage{enumitem} % [label=...] 옵션을 사용하기 위해 필수적인 패키지입니다.
\usepackage{setspace}
\graphicspath{{graph/}}
\geometry{a4paper, margin=1in}
\title{CLINIC 17-19}
\setstretch{1.3}

\begin{document}
\maketitle
\date{}

\begin{enumerate}
\item 다음 곡선과 직선 또는 곡선과 곡선으로 둘러싸인 도형의 넓이를 구하여라.
   % itemize 대신 enumerate를 사용하고, enumitem 패키지 덕분에 아래 옵션이 정상 작동합니다.
   \begin{enumerate}[label=(\arabic*)] 
    \item $y^2 = \dfrac{1-x}{x}$, $y=1, y=-1, x=0$
    \item $y=xe^{1-x}$, $y=x$
    \item $x=(y-2)^2+1$, $(x-1)^2+(y-1)^2=1$ ($x \ge 1$)
    \item $y^2 = 2x, \quad x^2 = 2y$
   \end{enumerate}

\item 좌표평면 위를 움직이는 점 P의 시각 $t$ 에서의 위치가 $(e^{-t} \cos t, e^{-t} \sin t)$ 일 때, 다음 물음에 답하여라.
   (1) 시각 $t$ 에서의 점 P의 속력을 구하여라.\\
   (2) $t=0$ 일 때부터 $t=a(a > 0)$ 일 때까지 점 P가 움직인 거리 $l$ 을 구하여라.\\

\item 함수 $f(x)=e^x+1$ 의 역함수를 $g(x)$ 라고 하자. $a$ 가 양의 상수일 때, 다음 정적분의 값을 구하여라.
\[ \int_0^a f(x) dx + \int_2^{f(a)} g(x) dx \]

\item 매개변수로 나타낸 함수 $x=a(\theta - \sin \theta), y=a(1 - \cos \theta)$ ($a>0, 0 \le \theta \le 2\pi$)의 그래프와 $x$ 축으로 둘러싸인 도형의 넓이를 구하여라.

\item $xy$평면에 곡선 $y = \dfrac{1}{x} \ln x$ 와 $x$축 및 직선 $x = e^2$ 으로 둘러싸인 도형이 있다. 이 도형을 밑면으로 하는 입체를 $x$축에 수직인 평면으로 자른 단면이 모두 반원일 때, 이 입체의 부피를 구하여라.

\item 곡선 $y = \sqrt{x}\sin x (0 \le x \le \pi)$ 와 $x$축으로 둘러싸인 도형을 $x$축 둘레로 회전시킨 입체의 부피를 구하여라.

\item $y = \sin x (0 \le x \le \pi/2), y=1, x=0$으로 둘러싸인 도형을 $y$축 둘레로 회전시킨 입체의 부피를 구하여라.

\item $y = e^x \sqrt{1-x^2}$과 $x$축으로 둘러싸인 도형을 $x$축 둘레로 회전시킨 입체의 부피를 구하여라.

\end{enumerate}

\end{document}